\chapter{Simulation 02: eight tetrahedral centers}\label{ch:s02}
\plain{This program moves eight centers through three coordinate directions.
The tetrahedra are fixed shapes drawn around them. Each center takes its turn
responding to an average of the currently stored center coordinates. The
movie shows the resulting shapes, but the missing numerical history limits
what can be measured from it.}

\section{The world and the drawing geometry}
The complete 119-line \filepath{3Dtriangles.py} initializes an $(8,3)$
position array uniformly through \filepath{np.random.rand(8,3)*10}
and an $(8,3)$ velocity array through
\filepath{(np.random.rand(8,3)-0.5)*0.2}. Thus positions lie in $[0,10)^3$
and velocity components in $[-0.1,0.1)$. No fixed seed is set.
Unlike S01, the initial velocities are generally nonzero.
There are 48 real-valued carried entries. \src{S02}{8--16}.
\claim{C30}

The relative vertex offsets are
\[
\tfrac12(1,1,1),\quad\tfrac12(-1,-1,1),\quad
\tfrac12(-1,1,-1),\quad\tfrac12(1,-1,-1),
\]
with face index triples $(0,1,2)$, $(0,1,3)$, $(0,2,3)$ and $(1,2,3)$.
Their edge length is $\sqrt2$, directly from the coordinate differences.
Every displayed tetrahedron retains its size and orientation. Its center
translates; its vertices do not have velocities, force laws or collision
rules. Face transparency is 0.6. No independent vertex wrapping is
implemented. \src{S02}{19--26,64--71,99--102}. \claim{C31}

\section{Exact centroid construction}
Write the componentwise position and displacement wraps as
\begin{equation}\label{eq:wrap10}
\W_{10}(z)=z\bmod10,\qquad
\D_{10}(z)=((z+5)\bmod10)-5.
\end{equation}
These map respectively to $[0,10)$ and $[-5,5)$ in each coordinate.
At a callback with supplied label $f$, let $x_i,v_i$ be the incoming
state and $x_i^+,v_i^+$ the resulting state. The source sets
\begin{equation}\label{eq:shared02}
t=0.1f,\qquad h=H(t),\qquad
g=\frac{1}{1+\frac18\sum_{j=0}^7\|v_j\|_2}.
\end{equation}
H and RPCO are the same formulas as
Equations~\eqref{eq:H} and~\eqref{eq:RPCO}. GRF accepts positions and
time as arguments but reads neither; it uses the global incoming
velocity array. It is computed once, not recalculated after each
body moves. \src{S02}{29--51,74--82}. \claim{C32}

For $i=0,\ldots,7$ in order,
\begin{align}
c_i&=\frac18\left(\sum_{j<i}x_j^++\sum_{j\ge i}x_j\right),\label{eq:centroid}\\
d_i&=\D_{10}(c_i-x_i),\qquad r_i=\|d_i\|_2,\label{eq:centroiddiff}\\
a_i&=0.1gh\,\frac{d_i}{r_i+0.01},\label{eq:kick02}\\
u_i&=v_i+a_i,\qquad v_i^+=\R(u_i),\label{eq:vel02}\\
x_i^+&=\W_{10}(x_i+0.1v_i^+).\label{eq:pos02}
\end{align}
The symbol $a_i$ denotes the source's velocity contribution, not a new
mechanism. The centroid includes the body itself and averages the
currently wrapped Cartesian coordinates. Earlier bodies have already
moved. The regularized direction has magnitude $r_i/(r_i+0.01)$,
multiplied by the common signed H and GRF factors. H is multiplied
exactly once. \src{S02}{83--96}. \claim{C33}

There is no binary switch in this program and no random draw after
initialization. Nevertheless, “attraction to centroid” is a source
comment rather than a guarantee of inward velocity. H can be negative,
and even a positive inward contribution may leave an outward carried
velocity after addition and damping. GRF lies in $(0,1]$ for finite
velocities and decreases as mean incoming speed grows.
\src{S02}{33--51,83--96}; \cite{construction02}. \claim{C34}

\section{Clock labels, animation and retained state}
The source requests 2,000 animation labels, a 50 ms interval and
\filepath{blit=True}. It supplies no separate initialization function.
It attempts a 60 fps MP4 save and then calls \filepath{plt.show()},
also showing after a caught save exception. Arrays are not reset between
save and display. All of this is at module level.
\src{S02}{107--119}. \claim{C35}

\begin{table}[htbp]\centering\small
\caption{Four distinct time-related quantities in S02. The first two affect
the numerical rule; the last two control requested presentation.}\label{tab:clocks02}
\begin{tabularx}{\linewidth}{@{}l Y@{}}\toprule
Quantity & Meaning\\\midrule
$\Delta t=0.1$ & Multiplier in both velocity addition and position advance\\
$t=0.1f$ & Argument passed to H on each callback; not an accumulator\\
GUI interval 50 ms & Nominal timer interval, not a measured callback duration\\
Movie 60 fps & Encoded presentation cadence, independent of model-step units\\\bottomrule
\end{tabularx}
\end{table}

If labels 0--1999 are processed once, H's argument ranges from 0 to
199.9, while 2,000 multipliers of 0.1 account for 200 model-step units.
That conditional count does not establish the historical number of
state-changing callbacks. The accepted local-library inspection
(Matplotlib 3.10.9) found initialization, save, interactive-start and
repeat paths that can invoke a callback with its first label. With
a mutating update function, an initialization draw is a numerical
update, not necessarily a harmless rendering of the initializer.
Historical dependency versions and backend events are unrecovered.
\cite{construction02}, TIME\_AND\_MOVEMENT and saved library excerpts.
\claim{C36}

Repeated labels can apply the same H value to changed arrays. A label
reset can renew the large starting gain while retaining positions and
velocities. This is a source/library possibility for S02; its movie
does not certify that exact callback chronology. S03's surviving logs
will supply direct evidence of repeated labels for that distinct
program. \claim{C37}

\section{Wrapping can affect both the image and the interaction}
From the position line, there is an integer coordinate vector $k_i$
such that
\begin{equation}\label{eq:seam}
x_i^+-x_i=0.1v_i^+-10k_i.
\end{equation}
The native increment $0.1v_i^+$ can therefore be small while a stored
coordinate difference is almost ten. The displayed shape moves with
the wrapped center; the movie does not draw a periodic seam connection
or interpolate a path across it. \src{S02}{95--101}. \claim{C38}

There is a second, distinct effect. Because $c_i$ is an arithmetic mean
of wrapped coordinates, a body's positive crossing contributes
$-10/8=-1.25$ to that coordinate of the group average, in addition
to ordinary motion. Later bodies in the same sweep see the changed
average. Wrapping $c_i-x_i$ afterward does not turn the preceding mean
into an intrinsic mean on a torus. The chosen displacement itself
also has a branch at a half-period. These are deductions from the
implemented operations, not proposed replacements.
\src{S02}{85--96}; \cite{construction02}. \claim{C39}

\plain{A boundary relocation can be a drawing jump and still matter to
later motion. This program averages the displayed coordinate
representatives before wrapping the direction. We must preserve that
choice to explain the historical construction.}

\section{The original movie and what it shows}
The attributed movie
\filepath{tetrahedral_attractors_animation.mp4} contains 2,000 decoded
frames at 60 fps, 1000 by 800 pixels, duration $33\frac13$ seconds.
Its name, title, geometry, axes and nominal encoding agree with the
source. The original lies in part\_1 while the source lies in part\_2;
the output's relative working-directory path makes that directory
difference insufficient to reject attribution.
No source/run receipt or ordered numerical history was recovered.
\cite{construction02}, VISUAL\_EVIDENCE. \claim{C40}

Figure~\ref{fig:F06} samples the full duration. Shapes become closer
in projection and later arrangements change much less visibly.
Several shapes remain separated in the final image. These observations
establish neither every three-dimensional pair distance nor exact rest.
The accepted study decoded all frames but visually reviewed selected
sheets and full-resolution stills; it did not claim to have manually
inspected every frame. \claim{C41}
\EvidenceFig{F06}

The consecutive window in Figure~\ref{fig:F07} locates a conspicuous
screen relocation between frames 13 and 14. It is consistent with the
source wrap operation, but the missing positions and velocities prevent
certification of the crossing axis and native displacement. No value
from S03's CSVs is substituted for this missing S02 history.
\claim{C42}
\EvidenceFig{F07}

Color is another observation variable. The source uses
\filepath{viridis((i+frame)\%8/8)}, so an individual body's color
changes even when its location changes little. The first eight frames
(Figure~\ref{fig:F08}) illustrate why color alone is not a stable
identity label. Projection and transparency can overlap tetrahedra
separated in depth; their silhouettes may look triangular despite
four-face source geometry. There are no accumulated trajectory lines
in this movie, unlike S01's static plot.
\src{S02}{64--71,99--102}. \claim{C43}
\EvidenceFig{F08}

The source supports coupling through the centroid, carried velocity
and shared speed feedback. The inspected movie supports visible
motion, one localized relocation and slower late changes.
It does not establish a leader, route-following pair, collision,
asymptotic equilibrium or physical gravitational interpretation.
These limits concern available evidence, not a theorem that the
program can never display another collective event. \claim{C44}
\FloatBarrier
